\begin{problem}[Warmup]{(a)}
    Explain why every normal subgroup is the kernel of some homomorphism.
\end{problem}
\unsolved

\begin{problem}[Warmup]{(b)}
    Explain why knowing the kernel of a homomorphism tells us essentially everything about the homomorphism.
\end{problem}
\unsolved

\begin{problem}[Warmup]{(c)}
    Can you tell if a homomorphism is onto, by just looking at its kernel?
\end{problem}
\unsolved

\begin{problem}[Warmup]{(d)}
    Can you tell if a homomorphism is one-to-one, by just looking at its kernel?
\end{problem}
\unsolved

\begin{problem}{1}
    Consider the function $\varphi : \mathbb{R} \to \mathbb{S}$ defined by
    \[\varphi(r) = e^{2\pi ri} = \cos(2\pi r) + i\sin(2\pi r),\]
    where $\mathbb{S}$ is the unit circle of Example 24.13. Show that $\varphi$ is an onto
    homomorphism. What is the kernel of $\varphi$? What two groups are isomorphic,
    according to the Fundamental Isomorphism Theorem 33.4?
\end{problem}
\unsolved

\begin{problem}{2}
    Consider the subgroup $G$ of $U(M_2(\mathbb{R}))$ given by
    \[G = \left\{\begin{pmatrix} a & 0 \\ b & 1 \end{pmatrix} : a \neq 0\right\}.\]
    (See Exercise 25.19.) Define $\varphi : G \to \mathbb{R}^*$ by letting
    \[\varphi\begin{pmatrix} a & 0 \\ b & 1 \end{pmatrix} = a.\]
    \begin{enumerate}
        \item[(a)] Prove that this is an onto group homomorphism.
        \item[(b)] What is the kernel of this homomorphism?
        \item[(c)] What two groups does the Fundamental Isomorphism Theorem assert are
        isomorphic?
        \item[(d)] What is the relationship between this exercise and Example 33.8?
    \end{enumerate}
\end{problem}
\unsolved

\begin{problem}{3}
    Suppose $H$ is a normal subgroup of $G$. Show that the natural homomorphism
    $\nu : G \to G/H$ is indeed a homomorphism, is onto, and has kernel equal to $H$, as
    claimed in Section 33.2.
\end{problem}
\unsolved

\begin{problem}{4}
    Let's redo Exercise 32.7, in light of the Fundamental Isomorphism Theorem. Suppose that
    $n$ is a positive integer and $d$ is a positive integer which (properly) divides $n$.
    Prove that
    \[\mathbb{Z}_n/\langle d \rangle\]
    is isomorphic to $\mathbb{Z}_d$.
\end{problem}
\unsolved

\begin{problem}{5}
    Complete the proof of the Fundamental Isomorphism Theorem 33.4, relying if necessary on
    the corresponding proofs for the Fundamental Isomorphism Theorem for Commutative
    Rings 19.1:
    \begin{enumerate}
        \item[(a)] Show that the function $\mu$ is well defined (that is, it does not depend
        on the coset representative chosen).
        \item[(b)] Show that the function $\mu$ is a group isomorphism (that is, that it
        preserves the group operation and is one-to-one and onto).
    \end{enumerate}
\end{problem}
\unsolved

\begin{problem}{6}
    Prove Theorem 33.2, using Theorem 17.2 as a model, if necessary: Suppose that
    $\varphi : G \to H$ is an onto group homomorphism, and $h \in H$. Prove that
    $\varphi^{-1}(h)$ is a coset of $\ker(\varphi)$.
\end{problem}
\unsolved

\begin{problem}{7}
    Let $G$ be a group and fix some $g \in G$. Consider the map $\varphi : \mathbb{Z} \to G$
    defined by $\varphi(n) = g^n$. Show that $\varphi$ is a group homomorphism. What are the
    possible kernels for $\varphi$? Describe the image of $\varphi$ in $G$.
\end{problem}
\unsolved

\begin{problem}{8}
    Consider the function $\varphi : \mathbb{C}^* \to \mathbb{R}^*$ defined by
    $\varphi(\alpha) = |\alpha|$. Prove that $\varphi$ is an onto homomorphism. What two
    groups are isomorphic, according to the Fundamental Isomorphism Theorem?
\end{problem}
\unsolved

\begin{problem}{9}
    Suppose that $G$ is a group with normal subgroups $H, K$. Consider the normal subgroups
    $H \cap K$ and $HK$. (See Exercises 32.8, 32.13, 32.14.) Prove that the groups
    $H/(H \cap K)$ and $HK/K$ are isomorphic. (You should compare this to the corresponding
    result in ring theory, in Exercise 19.23.)
\end{problem}
\unsolved

\begin{problem}{10}
    Consider the group $\mathbb{R}^3$, and suppose that $a, b, c \in \mathbb{R}$, not all
    zero. Define $\varphi : \mathbb{R}^3 \to \mathbb{R}$ by $\varphi(x, y, z) = ax + by + cz$.
    Show that $\varphi$ is a group homomorphism. What is the kernel of $\varphi$
    (considered \emph{geometrically})? What does the Fundamental Isomorphism Theorem say in
    this context? What is the geometric meaning of this assertion? (Compare this exercise
    to Exercise 28.9.)
\end{problem}
\unsolved

\begin{problem}{11}
    Consider the dihedral group $D_4$. Every element of $D_4$ can be written in the form
    $F^iR^j$, where $i = 0, 1$, and $j = 0, 1, 2, 3$, where $F$ is the `flip' about the
    vertical axis $(12)(34)$, and $R$ is the counterclockwise rotation $(1432)$. (See
    Exercise 28.12.) Define
    \[\psi : D_4 \to \mathbb{Z}_2 \times \mathbb{Z}_2\]
    by setting $\psi(F^iR^j) = ([i]_2, [j]_2)$.
    \begin{enumerate}
        \item[(a)] Prove that this is an onto homomorphism.
        \item[(b)] What is the kernel of this homomorphism?
        \item[(c)] The subgroup you obtained in part b is necessarily normal. (Why?) Prove
        this directly.
        \item[(d)] What two groups does the Fundamental Isomorphism Theorem assert are
        isomorphic?
    \end{enumerate}
\end{problem}
\unsolved

\begin{problem}{12}
    Suppose that $G$ is a group with normal subgroup $H$.
    \begin{enumerate}
        \item[(a)] Define a correspondence between the subgroups of $G$ containing $H$, and
        the subgroups of $G/H$, by assigning $K \to K/H$. Why is Exercise 32.14 relevant to
        this definition?
        \item[(b)] Show that the correspondence from part a is one-to-one and onto.
        \item[(c)] Furthermore, show that this correspondence also establishes a one-to-one
        correspondence between \emph{normal} subgroups of $G$ containing $H$ and
        \emph{normal} subgroups of $G/H$.
    \end{enumerate}
\end{problem}
\unsolved

\begin{problem}{13}
    Suppose that $H$ and $K$ are normal subgroups of the group $G$, and that
    $K \subseteq H$. Then the group $G/H$ is isomorphic to
    \[(G/K)/(H/K).\]
    (You should compare this to the corresponding result in ring theory, in
    Exercise 19.21.)
\end{problem}
\unsolved
